\documentclass[11pt,a4paper]{article}

\usepackage[margin=1in]{geometry}
\usepackage[utf8]{inputenc}
\usepackage[T1]{fontenc}
\usepackage{lmodern}
\usepackage{microtype}
\usepackage{xcolor}
\usepackage{hyperref}
\usepackage{enumitem}
\usepackage{titlesec}
\usepackage{tabularx}
\usepackage{booktabs}
\usepackage{fancyhdr}
\usepackage{parskip}

\definecolor{accent}{HTML}{d97706}
\definecolor{muted}{HTML}{6b7280}
\definecolor{codebg}{HTML}{f3f4f6}

\hypersetup{
  colorlinks=true,
  linkcolor=accent,
  urlcolor=accent,
  citecolor=accent,
  pdftitle={Missing Podo --- Build Report},
  pdfauthor={Sırrı Kaan Oğuzkan}
}

\titleformat{\section}
  {\Large\bfseries\color{accent}}
  {\thesection.}{0.6em}{}
\titleformat{\subsection}
  {\large\bfseries}
  {\thesubsection}{0.6em}{}

\setlist[itemize]{leftmargin=1.2em,itemsep=2pt,topsep=2pt}

\pagestyle{fancy}
\fancyhf{}
\fancyhead[L]{\footnotesize\color{muted}Missing Podo --- Build Report}
\fancyhead[R]{\footnotesize\color{muted}Jotform 2026 Frontend Challenge}
\fancyfoot[C]{\thepage}
\renewcommand{\headrulewidth}{0.3pt}

\newcommand{\code}[1]{\texttt{\small #1}}

\title{\vspace{-2em}\textbf{\Huge Missing Podo}\\[0.3em]\Large The Ankara Case --- Build Report}
\author{Sırrı Kaan Oğuzkan}
\date{April 2026}

\begin{document}

\maketitle
\thispagestyle{empty}

\vspace{-1em}
\begin{center}
\small\color{muted}
Jotform 2026 Frontend Challenge (Ankara) --- 3-hour timed build \\
Repository: \code{2026-frontend-challenge-ankara}
\end{center}

\vspace{1em}
\hrule
\vspace{1em}

\section{What was built}

An investigation dashboard that fetches submissions from five Jotform forms
(Checkins, Messages, Sightings, Personal Notes, Anonymous Tips), links records
across them by the people mentioned, and lets an investigator reconstruct
Podo's last movements.

\section{Stack and architecture}

\begin{itemize}
  \item \textbf{Client:} React 18 + TypeScript + Vite.
  \item \textbf{Server:} Node + Express + TypeScript.
  \item \textbf{Monorepo:} npm workspaces (\code{client/}, \code{server/}); single \code{npm run dev} boots both.
  \item \textbf{Data flow:} browser $\rightarrow$ \code{/api/records} on the server $\rightarrow$ Jotform API. The API key lives only in \code{server/.env}. The server normalizes the five forms into one shape \code{\{id, source, createdAt, timestamp, location, coordinates, people, text, urgency, confidence, fields\}} so the UI never sees raw \code{qid}s.
  \item \textbf{Key rotation:} \code{.env} accepts a comma-separated list; a 429 rotates to the next key automatically.
\end{itemize}

\section{Core features}

\renewcommand{\arraystretch}{1.35}
\begin{tabularx}{\textwidth}{@{}lX@{}}
\toprule
\textbf{Area} & \textbf{Summary} \\
\midrule
Fetch     & Parallel \code{Promise.all} over 5 form IDs; partial failures reported to UI instead of failing the whole page. \\
Normalize & \code{DD-MM-YYYY HH:mm} $\rightarrow$ ISO; \code{"lat,lng"} $\rightarrow$ \code{\{lat,lng\}}; skip \code{control\_head}/\code{control\_button}. \\
Link      & \code{groupByPerson} collects every name field (\code{person}, \code{partner}, \code{sender}, \code{recipient}, tip subjects), normalizes, returns people with their records. \\
Views     & List, Timeline, Map, Graph --- switchable, all share the same filtered set. \\
Search    & Matches across names, locations, and text. \\
Filter    & Source chips + time-range scrubber + selected person all compose. \\
Detail    & Full record view + "Other records mentioning these people" cross-links. \\
States    & Loading skeletons, error banner, partial-fetch warning, per-view empty states. \\
\bottomrule
\end{tabularx}

\section{Bonuses}

\subsection{Timeline view}
Records grouped by day, ordered by timestamp. Focuses on the selected person if any.

\subsection{Map view}
\begin{itemize}
  \item OSM tiles via react-leaflet.
  \item Markers colored per source; click opens a popup and selects the record.
  \item \textbf{Star pin} at Podo's last confirmed location, with a \textbf{1 km dashed magenta circle} for proximity.
  \item \textbf{Teal polyline} connecting Podo's geotagged appearances in chronological order.
  \item \textbf{Numbered step badges} on each unique Podo location (1 = earliest). Records at the same place share a number --- numbering reindexes against the current filter so it's always $1\dots N$ without gaps.
  \item \textbf{Legend overlay} explaining the star, steps, and source colors.
  \item Map $\leftrightarrow$ list sync: clicking a marker selects the record, selecting a record pans the map.
\end{itemize}

\subsection{Co-occurrence graph}
SVG radial layout. Podo centered, other people on the ring, edges weighted by how often they share records. Clicking a node selects that person.

\subsection{``Most suspicious'' scoring}
Heuristic surfaced in the Summary panel:
\begin{itemize}
  \item $+2$ per anonymous tip naming the person
  \item $+1$ per high-urgency message involving them
  \item $+1$ if they were seen within 1 km of Podo's last known location
\end{itemize}
Each entry expands to show which records drove the score, so the heuristic is legible instead of opaque.

\subsection{``Last seen with''}
Top co-occurrences for the selected person, drawn from Sightings and Checkins.

\subsection{Post-disappearance anomaly flag}
Records timestamped after Podo's last confirmed sighting get a ``$\triangle$~post-disappearance'' badge --- a quick way to spot leads or contradictions.

\subsection{Fuzzy name matching}
Turkish character folding (\code{ğ}$\rightarrow$\code{g}, \code{ş}$\rightarrow$\code{s}, etc.) + Levenshtein distance. Toggle in Controls so both modes are demoable. Selecting a person preserves selection across the toggle by re-mapping aliases.

\subsection{Time-range scrubber}
Dual-handle range over the data's min/max timestamps. Filters every view, including the graph and the map's numbered path.

\subsection{Shareable deep links}
URL hash serializes \code{view}, \code{person}, \code{q}, \code{sources}, \code{fuzzy}, \code{from}, \code{to}, \code{record}. A ``Share link'' button copies the current URL. Initial record deep-link is applied once on data load so closing the detail panel doesn't re-open it from stale hash.

\subsection{Keyboard shortcuts}
\begin{itemize}
  \item \code{/} focuses search
  \item \code{$\uparrow$} / \code{$\downarrow$} walks the visible record list
  \item \code{Esc} clears selection
\end{itemize}

\subsection{Theming}
Light/dark toggle backed by a \code{useTheme} hook. Persisted to \code{localStorage}, defaults to \code{prefers-color-scheme}. All colors routed through CSS custom properties (\code{--bg}, \code{--surface}, \code{--text}, \code{--text-soft}, \code{--muted}, \code{--accent}, \code{--chart-*}).

\subsection{Responsive pass}
Header stacks at $\leq 900$\,px, layout collapses to one column at $\leq 720$\,px, controls and sliders re-flow.

\section{Notable fixes along the way}

\begin{itemize}
  \item \textbf{Close-button race in the detail panel.} The deep-link effect was re-applying \code{hash.record} on every render, so closing the panel immediately re-opened it. Fixed by applying the initial deep link once via a ref.
  \item \textbf{Responsive header stayed right-aligned.} The base \code{.app\_\_header} rule was defined \emph{after} the \code{@media} block, so the cascade re-applied the non-responsive version inside the media range. Moved all media queries to the end of the stylesheet.
  \item \textbf{Map step numbers appeared to skip (1, 4, 8, 12\dots).} Multiple records were sharing coordinates, so markers stacked. Deduped steps by rounded lat/lng --- every unique place gets one step, and records at that place share its number.
  \item \textbf{\code{react-leaflet-cluster} incompatibility.} Clustering crashed at runtime against react-leaflet 4 + React 18, so it was dropped in favor of the numbered path + legend, which reads cleaner for 45 records anyway.
\end{itemize}

\section{Decisions and trade-offs}

\begin{itemize}
  \item \textbf{Server proxy over client-side fetch} --- required to keep the API key off the browser, and lets the server do the normalization once.
  \item \textbf{Client-side filtering/search} --- the dataset is bounded ($\sim 45$ records), so no server pagination or search index is needed.
  \item \textbf{CSS variables over a framework} --- kept the dep footprint tiny and made the theme toggle a two-line change.
  \item \textbf{Heuristic scoring, not ML} --- explainability matters more than accuracy; the breakdown under each score shows \emph{why}.
  \item \textbf{Numbered map path, not clustering} --- clearer for a small dataset; clustering was tried and dropped due to library incompatibility.
  \item \textbf{Commit + push per step} --- judging happens on commit history within the 3\,h window, so every verifiable checkpoint was committed.
\end{itemize}

\section{What's intentionally out of scope}

\begin{itemize}
  \item Server-side caching / rate-limit retries beyond simple key rotation.
  \item Automated tests.
  \item Backend pagination, database persistence, or write endpoints (forms are \code{DISABLED}).
  \item Marker clustering (attempted, dropped).
\end{itemize}

\section{How to run}

\begin{quote}
\code{npm install} \\
\code{cp server/.env.example server/.env}\quad\textit{\# set JOTFORM\_API\_KEY} \\
\code{npm run dev}\quad\textit{\# client + server together}
\end{quote}

Client at \url{http://localhost:5173}, API at \url{http://localhost:3001/api/records}.

\end{document}
